\section{Positive Transfer Along Alternatives (Option D)}\label{sec:option-d}

\subsection{Encoding structure}

Option D restricts the family to $(2,m)$.  Its exploratory encoder generalizes
the legacy two-series layout from four to five alternatives while keeping the
voter count fixed.  The variable basis contains direct series analogous to
\texttt{sel0} and \texttt{sel1}.  Bundled \texttt{minlib} emits both series;
the split package assigns one series to
\texttt{decisive\_voter0} and the other to
\texttt{decisive\_voter1}.

No voter-pair selector variables, pair-selection clauses, or all-voter split
are introduced.  Consequently, the bundled \texttt{minlib} clause multiset is
literally the union of the two split decisive-lever multisets over the same
variables.  The candidate renaming is therefore the identity:
\[
  \rho(v)=v.
\]
The checker nevertheless parses both clause multisets and re-verifies the map;
the conclusion does not rest on equal counts alone.

\subsection{Step~0 result}

Table~\ref{tab:option-d-step0} records the base-case control and the transferred
case.  Both packages are UNSAT in both cases, and the identity $\rho$ verifies
\CM.

\begin{table}[ht]
\centering
\caption{Option D clause-multiset comparison.  The final column decomposes the
bundled \texttt{minlib} clause count into the two split series.}\label{tab:option-d-step0}
\begin{tabular}{llll}
\toprule
Case & Classification & Status & Clause decomposition \\
\midrule
$(2,4)$ & \texttt{equiv\_cm\_persists} & UNSAT/UNSAT
  & $13{,}826=6{,}913+6{,}913$ \\
$(2,5)$ & \texttt{equiv\_cm\_persists} & UNSAT/UNSAT
  & $576{,}002=288{,}001+288{,}001$ \\
\bottomrule
\end{tabular}
\end{table}

The $(2,4)$ row controls the exploratory encoder against the known M1.5
structure.  The $(2,5)$ row is the positive M2.1 transfer result.  It shows that
increasing the alternative dimension does not force a change in the
clause-multiset relation when the two-voter witness basis is held fixed.

\subsection{Raw repair enumeration}

After review of the Step~0 result, raw lever-level repair enumeration was
authorized for Option D only.  The enumerator considers removal subsets of the
active bundled and split lever packages, solves candidates in increasing
cardinality order, and uses clause-removal monotonicity to avoid solving a
superset once a SAT repair is known.  Inclusion minimality is checked against
all strict subsets.

The repair families are the same at $(2,4)$ and $(2,5)$:
\begin{description}[style=nextline,leftmargin=2.2em]
  \item[Bundled.]
  $\lever{asymm},\ \lever{un},\ \lever{minlib},\
    \lever{no\_cycle4}$.
  \item[Split.]
  $\lever{asymm},\ \lever{un},\
    \lever{decisive\_voter0},\
    \lever{decisive\_voter1},\
    \lever{no\_cycle4}$.
  \item[Transported bundled.]
  $\lever{asymm},\ \lever{un},\
    \{\mathtt{decisive\_voter0},\mathtt{decisive\_voter1}\},\
    \lever{no\_cycle4}$.
\end{description}

The transported bundled family differs from the raw split family.  In
particular, \lever{minlib} becomes the two-lever set
\[
  \{\mathtt{decisive\_voter0},\mathtt{decisive\_voter1}\},
\]
whereas the split representation contains each decisive lever as an independent
singleton repair.  The classification is
\texttt{noncanonical\_persists} for both cases.

\subsection{Interpretation}

The $(2,5)$ result preserves both parts of the M1.5 witness:
\begin{enumerate}
  \item the bundled and split packages are \CM{} under identity $\rho$; and
  \item their raw repair presentations differ after transport by $\pi$.
\end{enumerate}
Because the underlying clauses are unchanged up to identity, the divergence is
not attributable to extra split clauses or encoding strength.  It arises from
how the same clause families are exposed as repair levers.

This is a positive transfer along one axis, not a general family theorem.
Keeping two voters is essential to the construction: there are exactly two
fixed witness series, so no choice among voter pairs is required.  The result
does not predict what happens when $n$ increases.
